\chapter{The Topology of the Gap}
\label{chap:topology-of-the-gap}

% Closed question: What can be said about the gap between recorded events
% without inventing structure not in the ledger?

\begin{chapteressay}

A ledger records discrete marks. Between two adjacent marks, the ledger
has no entries. What does that mean mathematically? Three possible
answers compete.

The first answer: the gap is empty. Nothing exists between the marks.
This answer is wrong. The gap between, say, the rational $1/2$ and the
rational $2/3$ contains the rational $7/12$, the irrational $\sqrt{2}/4
\cdot \pi/3$, and uncountably many other admissible values. The gap is
densely populated; it is not empty.

The second answer: the gap is whatever the recorder chooses to make it.
The recorder can fill the gap arbitrarily with any continuation that
fits the boundary marks. This answer is also wrong. The continuation
is constrained by structural conditions (continuity, smoothness,
periodicity) that the boundary marks impose. The recorder cannot
invent unrecorded structure; the recorder is constrained by what the
boundary forces.

The third answer, which is the chapter's: the gap has topology. The
gap is bounded by the surrounding marks. Within the bounds, the gap
admits a continuum of values constrained by the structural conditions
the record carries. The continuum is not blank; it has a specific
mathematical structure. The structure is what can be said about the
gap without inventing.

The chapter develops this through five mathematical structures. The
first is the density phenomenon: rationals and irrationals are both
dense in the reals; between any two marks lie infinitely many
admissible values, none of which the ledger records explicitly. The
second is admissible continuation: the Intermediate Value Theorem
and its generalizations force the existence of values within the
gap from the boundary data and the continuity condition. The third
is the metavariable: an unresolved value constrained by surrounding
marks but not yet assigned a specific number. The fourth is noise
and residue: the Gibbs phenomenon and Taylor remainders carry the
part of the original function that a chosen basis cannot represent.
The fifth is area as a gap quantity: a quantity determined by the
boundary alone, requiring no interior sampling.

The chapter's ground example is the rational/irrational dichotomy.
Both are dense in the reals; both are infinite; neither is empty;
yet they are disjoint. The density establishes that the gap between
two rationals contains continuum-many points. The chapter's claim
is that this topological structure is the natural domain of the
gap; the gap is not a blank space but a continuum bounded by the
marks.

The chapter's second concrete example is the damaged fresco. A
section of a fourteenth-century fresco has fallen away. The
restorer must continue the surviving paint into the lost area
without inventing structure. The 1948 archival photograph provides
external metavariable data: the lost area's appearance is recorded
in the photograph even though the wall no longer carries it. The
restorer's admissible continuation is constrained by the surviving
edges and by the photograph. The continuation is admissible only
because the structural constraints are recorded; otherwise the
restoration would be invention.

The chapter's closed question is:

\begin{center}
\emph{What can be said about the gap between recorded events without
inventing structure not in the ledger?}
\end{center}

The chapter's answer is: the gap has topology. The gap is densely
populated by admissible values; admissible continuations are forced
by structural constraints from the surrounding marks; metavariables
carry unresolved values whose specific values are constrained but
not yet determined; residues are the algebraic record of what an
approximation cannot capture; gap quantities (like area) are
determined by the boundary alone. None of these structures are
invented; all are forced by the surrounding marks and the
structural conditions the record carries.

The chapter's discipline is that the ledger records only what it
records. Inferences about the gap are admissible only when they
are constrained by the surrounding marks plus stated structural
conditions. Inferences that go beyond the constraints --- pure
invention --- are not admissible in the chapter's sense. The
distinction between admissible inference and invention is the
chapter's contribution to the book's discipline.

Subsequent chapters build on this. Chapter 4 examines how multiple
records combine through their boundaries. Chapter 5 introduces
distance from order: the gap's continuous structure becomes a
distance on the recorded marks. Chapter 6 introduces motion as the
selection of one continuous trajectory from many admissible ones.
Each chapter inherits the chapter's discipline: nothing is
admissible that is not constrained by the surrounding records.

\end{chapteressay}

\section{The Gap}
\label{se:the-gap}

Between any two distinct rational numbers, there is another rational
number. Between $1/2$ and $2/3$, for example, sits $7/12$. Between
$7/12$ and $2/3$ sits $5/8$. The process continues. The rationals are
\emph{dense} in themselves: any open interval of rationals contains
infinitely many rationals.

Between any two distinct rationals, there is also an irrational
number. Between $1$ and $2$ sits $\sqrt{2}$. Between $7/12$ and $2/3$
sits an irrational; specifically, any number of the form $7/12 +
\sqrt{2}/100$ lies in the interval and is irrational. The irrationals
are also dense in the rationals.

The two density statements together describe a peculiar structure.
The rationals are everywhere; the irrationals are everywhere; both are
infinite; neither is empty between any two distinct rationals. Yet the
rationals and the irrationals are disjoint. A rational is never an
irrational, and an irrational is never a rational.

\begin{phenom}{Density of Rationals and Irrationals}
\label{ph:density-rationals-irrationals}

\PhStatement Both the rationals $\mathbb{Q}$ and the irrationals
$\mathbb{R} \setminus \mathbb{Q}$ are dense in the real line: every
open interval contains infinitely many of each.

\PhOrigin The density of $\mathbb{Q}$ in $\mathbb{R}$ follows from the
Archimedean property of $\mathbb{R}$. The density of $\mathbb{R}
\setminus \mathbb{Q}$ in $\mathbb{R}$ follows from the uncountability
of $\mathbb{R}$ and the countability of $\mathbb{Q}$.

\PhObservation Between $0$ and $1$ there are countably infinitely
many rationals (e.g., $1/n$ for $n = 2, 3, 4, \ldots$) and
uncountably infinitely many irrationals. The irrationals are
strictly more numerous, yet the rationals are dense enough that
they leave no open subinterval untouched.

\PhConstraint The complement of $\mathbb{Q}$ in $\mathbb{R}$ is
not open: every irrational is a limit of rationals. The complement
is not closed either: every rational is a limit of irrationals.
The rationals are neither open nor closed in $\mathbb{R}$.

\PhConsequence The gap between two consecutive entries in any
rational ledger is densely filled by irrationals, none of which the
ledger records. The gap has structure: it carries a continuum of
points that the ledger cannot represent.

\end{phenom}

The density phenomenon establishes the chapter's central observation:
the gap between two recorded marks is not empty. The gap contains
infinitely many admissible values that the ledger does not record. The
gap is, in this precise sense, a topological space in its own right.

The chapter's first claim is that the gap is admissible. By default,
the ledger records discrete marks. Between marks, the ledger has no
entries. But the gap is not blank in the mathematical sense; it is
densely populated by potential entries that the ledger could in
principle accept, but has not. The gap is the topological structure
the ledger carries implicitly.

What can be said about the gap without inventing structure? The
question is the chapter's closed question. The answer must respect
the discreteness of the ledger while admitting that the gap has
internal structure. The gap is a topological object; the ledger is
a discrete object; the relationship between them is the chapter's
subject.

The simplest mathematical answer is: the gap is bounded above and
below by the surrounding marks. Between the rational marks $1/2$ and
$2/3$, the gap is the open interval $(1/2, 2/3)$. The interval has
a length: $2/3 - 1/2 = 1/6$. The length is the gap's measure.

Below the surface, more can be said. The interval $(1/2, 2/3)$
contains $7/12$ (a rational) and $\sqrt{2}/\pi$ (an irrational,
approximately $0.45$, not in the interval --- example unreliable;
better: $1/2 + (\pi - 3)/8 \approx 0.518$). The interval contains
countably many rationals and uncountably many irrationals. The
cardinality of the interval is $2^{\aleph_0}$ (continuum). All of
this is mathematically present without being recorded in the ledger.

The chapter's discipline is that none of this structure is invented;
all of it is derived from the surrounding marks and the order's
properties (density). The structure is forced by the partial order
on the rational ledger together with the assumption that the order
admits Cauchy completion.

\section{Admissible Continuation}
\label{se:admissible-continuation}

The Intermediate Value Theorem says: if $f$ is a continuous function on
a closed interval $[a, b]$, with $f(a) < 0$ and $f(b) > 0$, then there
exists a point $c \in (a, b)$ where $f(c) = 0$. The theorem is a
statement about the gap: between two recorded values with opposite
signs, the function must cross zero somewhere; the zero is in the
gap; the gap cannot be empty in the sense of admitting no zero.

The theorem's proof uses the completeness of the real numbers. Define
$S = \{x \in [a, b] : f(x) < 0\}$. $S$ is non-empty ($a \in S$) and
bounded above ($b$ is an upper bound). By completeness, $S$ has a
supremum $c \in [a, b]$. Continuity of $f$ implies $f(c) = 0$: if
$f(c) < 0$, then $f$ is negative in a neighborhood of $c$, contradicting
the supremum property; if $f(c) > 0$, similarly.

The theorem is the mathematical statement that the gap admits an
admissible continuation. Given the boundary values $f(a) < 0$ and
$f(b) > 0$, the gap between them must contain a zero. The continuation
exists; it is forced by the boundary data and the continuity. The
specific zero may not be explicitly computed by the ledger, but its
existence is admissible.

The chapter's second claim is that admissible continuations are
mathematical objects that exist independently of the ledger's
ability to record them. The continuation may be specified
implicitly by boundary data and structural constraints (continuity,
smoothness, periodicity). The continuation may not be a single
explicit value; it may be a parameterized family that the boundary
data constrains.

The Intermediate Value Theorem is the cleanest example. The
boundary data are $f(a) < 0$ and $f(b) > 0$. The structural
constraint is continuity. The admissible continuation is the
existence of a zero in $(a, b)$. The continuation is admissible
even though the ledger does not record the zero's value
explicitly.

More elaborate examples follow. The Mean Value Theorem says: if
$f$ is continuous on $[a, b]$ and differentiable on $(a, b)$, then
there exists $c \in (a, b)$ with $f'(c) = (f(b) - f(a))/(b - a)$.
The theorem is admissible because the boundary data ($f(a)$,
$f(b)$) and the structural constraints (continuity,
differentiability) together force the existence of $c$. The
ledger may not record $c$ explicitly; the ledger admits its
existence as a constrained metavariable.

In Lean, the analog appears in the typeclass cascade for
admissible continuations. \texttt{NoisyProcess} certifies that a
process produces an output even when intermediate values are not
directly observable. \texttt{PHYSICAL} certifies that the
process is grounded in admissible measurement. Together they
allow the ledger to assert the existence of an output without
naming its specific value.

The Lean construction is precise. Suppose a function is defined
piecewise: $f(x) = -1$ for $x < 0$, $f(x) = 1$ for $x > 0$. The
Intermediate Value Theorem does not apply: the function is not
continuous at $x = 0$, so there is no admissible zero in the
gap. The ledger admits this: $f$ is not \texttt{REPRESENTABLE} as
a continuous function. Without the representable certificate, no
admissible continuation is guaranteed.

The chapter's lesson is that admissibility is a structural
constraint, not an arbitrary one. The Intermediate Value Theorem
requires continuity; without continuity, the theorem fails and
the gap may contain no admissible continuation. With continuity,
the gap is constrained to admit the continuation. The ledger's
discipline is to record the constraints; the constraints determine
the admissibility of continuations.

\section{The Metavariable}
\label{se:the-metavariable}

The quadratic formula
\[
x = \frac{-b \pm \sqrt{b^2 - 4ac}}{2a}
\]
solves $ax^2 + bx + c = 0$ for $x$. The formula is admissible as a
mathematical statement before any specific values for $a, b, c$ are
given. The symbols $a, b, c$ are \emph{metavariables}: placeholders
whose values have not been resolved, but which carry structural
constraints (they appear in specific syntactic positions and are
related by specific arithmetic).

The formula is mathematical content even unresolved. The solutions
have certain properties that follow from the formula's structure
regardless of the specific values:
\begin{enumerate}
\item The two solutions sum to $-b/a$ (Vieta's formula).
\item The two solutions multiply to $c/a$ (Vieta's formula).
\item The solutions are real when the discriminant $b^2 - 4ac \ge 0$.
\item The solutions are complex conjugates when the discriminant
$b^2 - 4ac < 0$.
\end{enumerate}
All four statements are admissible in the unresolved form. They
constrain what the resolved values can be without naming them
specifically.

This is the chapter's third claim. A metavariable is an admissible
unresolved value: a placeholder constrained by the structural
relations of the ledger but not yet assigned a specific value.
The metavariable is part of the ledger's record; it is carried
honestly as ``unresolved'' rather than being silently filled with a
default value.

The metavariable concept is essential for the chapter's question
about the gap. The gap between two recorded marks may contain
admissible values, but the ledger does not record them explicitly.
The admissible values are metavariables: their existence is
asserted by the surrounding marks' structural constraints; their
specific values are unresolved.

In Lean, the metavariable is a fundamental construct. The
elaborator routinely creates metavariables during type inference:
when the user writes \texttt{def x := f y} without specifying
the return type, the elaborator creates a metavariable for the
type, then resolves it from the context. The metavariable is
present in the elaboration state until the constraints from
the surrounding code determine its value.

\begin{leanexcerpt}{Metavariable}
inductive Metavariable ($\alpha$ : Type i) : Type (i + 1) where
  | base : Fact $\to$ $\alpha$ $\to$ Metavariable $\alpha$
  | step : Fact $\to$ Metavariable $\alpha$ $\to$ Metavariable $\alpha$
\end{leanexcerpt}

A \texttt{Metavariable} in the project's Lean code is an inductive
type over a carrier $\alpha$. Each layer carries a \texttt{Fact}
(the bookkeeping that something has been refined) and either a
witness value at the base or another metavariable at the step. The
inductive layers correspond to successive bisections of the gap:
each \texttt{step} sharpens what was given at \texttt{base} without
ever resolving to a concrete value.

The \texttt{COMPARABLE} typeclass in \texttt{Episode3.lean}
formalizes the comparison of metavariables. Two metavariables are
\texttt{COMPARABLE} when their types are consistent: that is, when
they could resolve to compatible values. The comparison is
structural; the resolution may not have occurred yet.

The chapter's claim is that the metavariable is the algebraic
shadow of the gap. The gap contains values; the ledger does not
record them; the metavariable represents the unresolved entries.
The metavariable's constraints are inherited from the surrounding
records; the resolution of the metavariable is the algebraic
event of inserting a specific value into the gap.

Two examples sharpen the concept.

First, the quadratic formula. The discriminant $b^2 - 4ac$ is a
metavariable: it is computable from $a, b, c$ when those are
resolved, but the formula carries it before resolution. The
behavior of the formula (real or complex roots) depends on the
sign of the discriminant. The sign is a property of the
metavariable that the ledger can track without resolving the
specific values.

Second, the function $f(x) = \sin(x)/x$ at $x = 0$. The function
is not defined at $x = 0$ by the explicit formula. The limit as
$x \to 0$ is $1$ (by L'Hopital's rule or Taylor series). The
admissible continuation of $f$ to $x = 0$ is $f(0) = 1$. Before
the limit is taken, $f(0)$ is a metavariable: known to exist by
the surrounding structure, not yet resolved by the explicit
formula. After the limit is taken, the metavariable resolves to
$1$. The ledger records the metavariable; the resolution adds the
specific value.

The chapter's discipline: every admissible continuation that the
ledger does not record explicitly is a metavariable. The
metavariable carries the structural constraints that the gap
imposes. The resolution of the metavariable is the algebraic
event of supplying a specific value. The unresolved metavariable
is honest; the resolved value is admissible.

\section{Noise and the Admissible Continuation}
\label{se:noise-and-the-admissible-continuation}

The square wave function $s(x)$ takes value $+1$ for $0 < x < \pi$ and
$-1$ for $\pi < x < 2\pi$, extended periodically with period $2\pi$.
Its Fourier series is
\[
s(x) = \frac{4}{\pi}\left(\sin x + \frac{\sin 3x}{3} + \frac{\sin 5x}{5}
+ \frac{\sin 7x}{7} + \cdots\right).
\]
The series is an infinite sum of smooth functions. Each partial sum
$s_N(x)$ uses only the first $N$ terms; each $s_N$ is a smooth
function. The infinite sum converges to $s$ at every point of
continuity, but near the discontinuity at $x = 0$ (or $x = \pi$),
the convergence is non-uniform.

Near the discontinuity, the partial sums overshoot. At every fixed
$N$, the partial sum $s_N(x)$ achieves a maximum of about
$1.1789 = 1 + 2\int_0^\pi \frac{\sin u}{u} du / \pi - 1$ just before
the discontinuity, regardless of how large $N$ is. The overshoot does
not disappear as $N \to \infty$; it just narrows: the overshoot's
width shrinks to zero, but its height stays at about $1.179$.

This is the Gibbs phenomenon. It is the chapter's instance of an
admissible continuation that carries a residue. The smooth basis
(sines and cosines) cannot exactly represent the discontinuity;
every finite partial sum has the overshoot. The overshoot is the
mathematical record of the gap between the discrete (the
discontinuous square wave) and the smooth (the basis functions).

The Gibbs phenomenon is not noise in the statistical sense. It is
deterministic: the same square wave produces the same Gibbs
overshoot every time the Fourier series is computed. The
overshoot is a structural feature of the approximation. It is the
residue of the admissible continuation by smooth functions.

The chapter's fourth claim is that admissible continuations carry
residues. The smooth approximation captures part of the original;
the residue is what the approximation cannot capture. For a
discontinuous function approximated by smooth functions, the
residue is the Gibbs overshoot. For a non-smooth function (one
with a cusp or corner), the residue is concentrated at the non-
smooth point.

The general principle: an admissible continuation by a basis $B$
of a function $f$ produces an approximation $f_B$ and a residue
$r = f - f_B$. The residue is the part of $f$ that $B$ cannot
represent. The residue may be zero (perfect approximation) or
non-zero (imperfect approximation). The structure of the residue
characterizes the basis's limitations.

In Lean, the residue appears as the \texttt{RESIDUE} typeclass
introduced in Chapter 1. The class certifies the existence of a
remainder after admissible decomposition. For the Gibbs phenomenon,
the \texttt{RESIDUE} is the overshoot; the \texttt{quotient} is
the partial sum; the \texttt{reconstruct} is the proof that
partial sum plus overshoot equals the original square wave (in
the appropriate limit sense).

A second example. The Taylor series of $e^x$ around $x = 0$ is
\[
e^x = 1 + x + \frac{x^2}{2!} + \frac{x^3}{3!} + \cdots.
\]
The series converges for all real $x$, but only the infinite sum
equals $e^x$ exactly. Any finite partial sum is an approximation
with residue. For $x = 1$, the partial sum to $n$ terms equals
$e - r_n$ where $r_n$ is bounded by $1/n!$. The residue shrinks
factorially.

The Taylor residue is well-behaved: it decreases monotonically
with the number of terms. The Gibbs residue is not: the overshoot
does not decrease with the number of terms. The difference is in
the smoothness of the original function. $e^x$ is analytic
(infinitely differentiable everywhere); the square wave is
piecewise constant with jumps. The smooth function admits a
clean Taylor approximation; the discontinuous function admits
only the Gibbs-suffering Fourier approximation.

The chapter's lesson is that the residue's structure reflects the
gap's structure. A clean residue (vanishing with the basis size)
indicates that the basis is well-matched to the function. A
persistent residue (not vanishing with the basis size) indicates
a structural mismatch. The mismatch is the admissible record of
what the basis cannot capture.

The discipline of admissible continuation is therefore to carry
the residue explicitly. An approximation that pretends to capture
the original without residue is overstating its accuracy. An
approximation that explicitly carries the residue is honest. The
ledger admits both the approximation and the residue; together
they reconstruct the original.

\section{Area as a Gap Quantity}
\label{se:area-as-a-gap-quantity}

Consider a closed curve in the plane: a circle, an ellipse, a heart-
shaped curve, a polygon. The curve bounds a region of the plane. The
region has an area: a non-negative real number that measures its
extent. The area is a gap quantity in the chapter's sense: it depends
only on the curve's bounding shape, not on what is inside or how the
inside is traversed.

The area of a region $R$ bounded by a closed curve $\gamma$ can be
computed by Green's theorem:
\[
\text{Area}(R) = \frac{1}{2} \oint_\gamma (x \, dy - y \, dx).
\]
The integral is over the boundary curve, not over the interior. The
boundary determines the area; the interior need not be explicitly
sampled. The area is the gap quantity: it lives in the interior
without being sampled at any interior point.

The chapter's fifth claim is that gap quantities are quantities
determined by the boundary alone. Green's theorem is the cleanest
example: the area is a property of the interior, but the integral
formula expresses it purely as a boundary integral. The interior's
contents are unnecessary for computing the area; only the boundary
is needed.

This is a profound statement about the chapter's question. Between
two recorded marks (the boundary), the gap contains structure (the
interior). Some quantities depend on the interior's specific
contents (these are not gap quantities). Some quantities depend only
on the boundary (these are gap quantities). The chapter's
admissible continuations focus on the latter.

The area is path-independent in the sense that two parameterizations
of the same closed curve produce the same area. Reparameterizing the
curve (traversing it faster or slower, starting at a different
point) does not change the area. The area is intrinsic to the
curve's shape.

In Lean, the gap quantity concept appears as the typeclass
\texttt{Area} in \texttt{Episode3.lean}. The class certifies the
existence of an admissible area for a closed region:

\begin{leanexcerpt}{Area}
inductive Area
  | t : Fact $\to$ Area
  | dt : Fact $\to$ Number $\to$ Area $\to$ Area
\end{leanexcerpt}

The inductive carries the gap quantity by accumulation. The
\texttt{t} constructor records the initial \texttt{Fact} (a base
slice of admissible witness). The \texttt{dt} constructor records
an increment: a \texttt{Fact}, a \texttt{Number} contribution, and
the prior \texttt{Area} on which the contribution is layered. The
total quantity is what the accumulated stack reports, not a single
number measured by a function on the boundary.

The Stokes theorem generalizes Green's theorem to higher
dimensions. For a vector field $F$ on a manifold $M$ with boundary
$\partial M$, the integral of $F$ over $\partial M$ equals the
integral of $dF$ (the exterior derivative) over $M$:
\[
\int_{\partial M} F = \int_M dF.
\]
The left side is a boundary integral; the right side is an interior
integral. The two sides are equal: the boundary determines the
interior's contribution, and vice versa. The relationship is the
universal version of the chapter's gap-quantity concept.

Examples of gap quantities include:
\begin{itemize}
\item Area enclosed by a planar curve (Green's theorem).
\item Volume enclosed by a closed surface (divergence theorem).
\item Flux of a vector field through a closed surface.
\item Total charge in a region (Gauss's law).
\item Winding number of a closed curve around a point.
\end{itemize}
Each is determined by the boundary alone; each is a gap quantity in
the chapter's sense.

The chapter closes with this observation. The gap between recorded
marks may carry quantities that depend only on the boundary. These
quantities are admissible in the ledger even though the interior is
not directly recorded. The boundary's structure determines the gap
quantity; the interior's specific contents do not. The admissible
continuation by gap quantities is the most economical: it carries
the boundary record and computes the interior content from it.

\begin{bridge}

The chapter's question was what can be said about the gap between
recorded events without inventing structure not in the ledger.

The answer is: the gap has topology. The rationals are dense; the
irrationals are dense; between any two rational marks there are
infinitely many admissible values neither recorded. Admissible
continuations are forced by structural constraints (continuity,
smoothness); the Intermediate Value Theorem is the canonical
example. Metavariables are unresolved values whose constraints
are inherited from surrounding marks; they admit eventual
resolution by additional data. Noise and approximation residues
(Gibbs phenomenon, Taylor remainder) carry the part of the
function the chosen basis cannot capture. Area is a gap quantity:
determined by the boundary alone, requiring no interior sampling.

What remains is composition. The chapter has identified the
topological structure of a single gap. Multiple records combine
through their boundaries; the chapter has not yet examined how.
Chapter 4 takes up the algebra of records: how multiple records
combine into a structured composite without losing their
individual identities.

\end{bridge}

\begin{coda}{The Damaged Fresco}

The fresco occupies the south wall of a small chapel in the hills.
The chapel was painted in the late fourteenth century by a journeyman
artist whose name is uncertain. The fresco shows a scene from the
life of a saint: the figures arranged in a procession, a landscape
behind them, the saint at the center with arms raised in blessing.
The wall is approximately three meters wide and two meters high.

Time has not been kind. A section of plaster about thirty centimeters
square has fallen away from the central area, taking with it most of
the saint's torso and the upper half of two attendant figures.
Smaller losses dot the rest of the wall. The colors have darkened
unevenly: some pigments have proven more stable than others. The
overall image is recognizable, but specific details in the damaged
area cannot be read from what remains.

The restorer arrives with a portfolio of the chapel's archival
photographs: a black-and-white image from 1923, a color image from
1948, and a high-resolution digital scan from 2007. The 1948 image
shows the fresco in better condition; the central damage had not yet
occurred. The 2007 scan shows the wall in roughly its current state.

The restorer's task is to assess what can be inferred about the
missing paint and, where possible, to restore the surface so that
the fresco's content is again legible. The work proceeds slowly. The
restorer cannot invent the missing paint; she can only continue the
surviving structure into the gap.

She begins with the boundary of the loss. The plaster's edges are
sharp where the fall has been recent (or where the conservation team
has cleaned the edge); they are softer where older damage has
weathered. The intact paint along the boundary tells her what
colors were present at the edge: the saint's robe was a dark blue
with a red trim; the attendants' garments were earth tones; the
landscape behind was a muted green.

These boundary colors are the recorded marks. The lost area is the
gap. The chapter's question --- what can be said about the gap
without inventing structure --- is the restorer's question.

She consults the 1948 image. The image shows the central area
intact: the saint's torso visible, the attendants' upper bodies
visible, the brushwork's character readable. The 1948 image is
not a recorded mark in the fresco itself; it is an external
record that can be compared to the surviving paint. The
comparison is admissible: the 1948 image and the surviving paint
should agree where both are present.

The comparison succeeds at most points. The 1948 image shows the
robe's color matching the surviving edge. The image shows the
attendants' garments matching the surviving edge. The image shows
the landscape green matching the surviving edge. The agreement is
the calibration: the 1948 image is a reliable record of the
fresco's state at that time.

For the lost area, the 1948 image shows: the saint's right hand
extended in blessing, the robe's drape over the right shoulder,
the upper body of the right attendant. These details are not
recorded in the surviving paint --- they are gone --- but they are
recorded in the 1948 image. The image is an external mark that
constrains the gap.

The restorer can therefore propose an admissible continuation of
the surviving paint into the lost area, guided by the 1948 image.
The continuation is bounded by the surviving edges and by the
image's recorded content. The continuation is not free invention;
it is constrained by external evidence.

She does not, however, paint directly into the lost area. The
conservation discipline is to mark the restored portions visibly:
either by leaving them unpainted (so the lost area is obvious as
loss), or by painting with a different technique (so the
restoration is distinguishable from the original on close
inspection). The choice depends on the institution's
conservation philosophy. For this chapel, the choice has been
made: the lost area will be filled with a neutral tone matching
the wall behind, with the figures' outlines suggested by careful
tratteggio (a technique of parallel fine lines that produces an
overall color from a distance but is identifiable as restoration
on close view).

The work proceeds. The restorer prepares the wall surface (cleaning,
consolidating the surrounding plaster, applying a fresh ground
layer in the lost area). She mixes pigments to match the
surrounding colors, accounting for the original's likely fading
under four centuries of light. She paints the suggested outlines
using tratteggio. The result is a wall that is again readable as
a scene: the saint's torso, the attendants' figures, the
landscape behind --- all visible at viewing distance, with the
restoration distinguishable on close inspection.

The chapter's structure appears throughout. The lost area is the
gap. The surviving paint is the boundary. The 1948 image is the
external metavariable: a recorded value of what the gap should
contain, even though the wall itself does not currently record
it. The continuation is admissible because it is constrained
by the boundary and the metavariable; it is not free invention.

What if the 1948 image had not existed? The restorer would have a
more difficult task. She could still continue the surviving
edges (the robe's blue, the landscape's green) into the lost
area. She could not, without external evidence, determine the
specific posture of the saint's torso or the attendants'
positions. The admissible continuation would be more cautious:
filling the gap with the boundary colors but not asserting
specific figures within it. The image's absence would mean the
gap's interior structure is unresolved.

This is the chapter's discipline. The restorer continues what is
forced by the boundary and the constraints. The restorer does
not invent unrecorded structure. The 1948 image is recorded
structure; using it is admissible. The restorer's own
imagination is not recorded structure; using it as a source of
content would be inadmissible.

The fresco's area, in the technical sense of Green's theorem, is
determined by the wall's perimeter. The total area is roughly
six square meters. The restored area is roughly thirty by thirty
centimeters: nine hundredths of a square meter. The percentage
of restoration is about one and a half percent of the total
fresco. The percentage is documented in the conservation
report.

The Gibbs phenomenon analog appears too. Where the surviving
paint meets the restoration, the transition is sharp on close
inspection. The restorer's tratteggio is visibly distinct from
the original brushwork. The transition is the residue of the
restoration: the part of the work that cannot be made invisible
without falsifying the record. The discipline is to leave the
transition visible so that the restoration is identifiable.

At the end of the project, the chapel's south wall is again
legible. A visitor at viewing distance sees the saint, the
attendants, and the landscape. A visitor on close inspection sees
the restoration's visible marks. The fresco is restored honestly:
not as it was originally, but as a continued record of what has
survived plus what the external evidence constrains.

The chapter's algorithmic structure is the restorer's procedure.
The gap is the lost area. The boundary is the surviving paint.
The metavariable is the 1948 image. The admissible continuation
is the tratteggio-filled outline. The residue is the visible
transition. None of the operations invent unrecorded structure;
all are constrained by the surrounding marks and external evidence.

\end{coda}
